\documentclass[10pt,a4paper]{amsart}
\usepackage[margin=19mm]{geometry}
\usepackage{amsmath,amssymb,mathtools}
\usepackage[scr=rsfs,cal=euler]{mathalfa}
\usepackage{xfrac}
\usepackage[unicode,hidelinks,bookmarks=false]{hyperref}
\newcommand{\ie}{i.e.\ }
\newcommand{\cf}{cf.\ }
\newcommand{\resp}{respectively\ }
\newcommand{\Subsection}{\subsection}
\providecommand{\nobreakdash}{\nobreak\hbox{-}}
\setlength{\emergencystretch}{2em}
\multlinegap0pt
\usepackage{amssymb}
\usepackage{bm}
\usepackage{tikz}
\usepackage{float}
\newtheorem{theorem}{Theorem}
\newtheorem{lemma}{Lemma}
\def\NN{\mathbb{N}}
\def\SS{\mathbb{S}}
\def\mP{\mathcal{P}}
\def\mq{\mathfrak{q}}
\def\pc{\mathcal{P}^c}
\def\sul{\sum\limits}
\title{Optimal Neural Network Approximation via Empirical Least Squares with Deterministic Samples}
\author{Xinliang Liu, Tong Mao, Jinchao Xu}
\date{Source: arXiv:2608.06687v1 (CC BY 4.0)}
\begin{document}
\maketitle

\noindent\emph{Source and licence.} Optimal Neural Network Approximation via Empirical Least Squares with Deterministic Samples. Version arXiv:2608.06687v1 (CC BY 4.0), distributed by arXiv under the Creative Commons Attribution 4.0 International licence, http://creativecommons.org/licenses/by/4.0/, as recorded in the arXiv licence metadata for this exact version and shown on the abstract page https://arxiv.org/abs/2608.06687v1. Full licence notice: \texttt{licenses/arxiv-2608.06687-CC-BY-4.0.txt}.\par\smallskip

\noindent\textit{Editorial scope.} Counted unit: the article's theorem thm:main\_sphere (source0.tex lines 1152--1162). The article's complete original proof of it is delivered with it (source0.tex lines 1344--1541, one \texttt{proof} environment of 198 lines, which the article places away from the statement). No mathematics is written, repaired or re-derived: the statement and the proof are the article's own lines, in the article's own notation, and no retained line is substituted or re-flowed.  Retained uncounted as the source-local context the proof needs: lines 162--164; lines 209--211; lines 411--413; lines 415--417; lines 540--557; lines 553--555; lines 675--685; lines 682--684; lines 1089--1106; lines 1168--1173; lines 1241--1251; lines 1305--1310.  Nothing was omitted from any retained span, so the package carries no omission record.  No retained line contains a citation, so the wrapper carries no bibliography and no bibliography entry.  No retained line contains a figure environment, an image-inclusion command or drawing code, so no image is delivered and the package declares no asset dependency.  Editorial formatting only, in addition: an AMS \texttt{amsart} preamble that restates the article's own declarations and macros used by the retained lines, namely the environments \texttt{theorem}, \texttt{lemma} on separate counters, together with the standard packages those lines need.  The retained environments print in this document as Theorem 1, Lemma 1 in order of appearance, whereas the article prints the same items with the numbering of its own full document.  The complete native source of the article as distributed in the arXiv e-print for this exact version is bundled unchanged as \texttt{sources/207023/source0.tex}.
\begin{equation}\label{eqn:f_evenodd_relu}
f(\eta)=(-1)^{k+1}f(-\eta),\qquad\eta\in\SS^d,
\end{equation}

\begin{equation}\label{eqn_def_collocation}
u_{n,m}\in\arg\min_{v_n\in L_n^k(\Theta_n)}\frac1m\sum_{i=1}^m\left(f(\eta_i^*)-\mathfrak L_\beta v_n(\eta_i^*)\right)^2.
\end{equation}

\begin{equation}\label{eqn:def_operator}
\mathfrak L_\beta v=\sul_{m=0}^\infty\ell_\beta(m)\Pi_m v,
\end{equation}

\begin{equation}\label{eqn:elliptic_symbol}
c_{\mathfrak L}(1+m)^\beta\leq\ell_\beta(m)\leq C_{\mathfrak L}(1+m)^\beta,\qquad m\geq0.
\end{equation}

\begin{theorem}\label{thm:inverse}
    Let $d,n\geq2$, $k\in\NN_0$, $0\leq s<r< k+\frac{1}{2}$, and let $\Theta_n=\{\theta_{j,n}^*\}_{j=1}^n\subset\SS^d$ be fixed with antipodal separation $\underline h>0$. With the simplified notation $\theta_j^*=\theta_{j,n}^*$, for any $f_n\in L_n^k(\Theta_n)$ the following Bernstein inequality holds:
    \begin{equation}
        \|f_n\|_{\mathcal{H}^r(\SS^d)}\lesssim \underline{h}^{-(r-s)}\|f_n\|_{\mathcal{H}^s(\SS^d)}.
    \end{equation}
    
    In particular, if $\Theta=\{\Theta_n\}_{n\ge1}$ is an antipodally quasi-uniform sequence,
    \begin{equation}
        \|f_n\|_{\mathcal{H}^r(\SS^d)}\lesssim n^{\frac{r-s}{d}}\|f_n\|_{\mathcal{H}^s(\SS^d)},
    \end{equation}
    where the implied constant is independent of $n$, $\Theta_n$, and $f_n$.

    Moreover, let $\beta\geq0$ and $\mathfrak L_\beta$ satisfy \eqref{eqn:def_operator}--\eqref{eqn:elliptic_symbol}. For $0\leq s<r< k+\frac{1}{2}-\beta$, and any $v_n\in\mathfrak{L}_\beta L_n^k(\Theta_n)$,
    \begin{equation}\label{eqn:residual_bernstein}
\|v_n\|_{\mathcal H^r(\SS^d)}\lesssim\underline h^{-(r-s)}\|v_n\|_{\mathcal H^s(\SS^d)}.
\end{equation}

\end{theorem}

    \begin{equation}
\|v_n\|_{\mathcal H^r(\SS^d)}\lesssim\underline h^{-(r-s)}\|v_n\|_{\mathcal H^s(\SS^d)}.
\end{equation}

\begin{theorem}\label{thm_inverse_hig_freq}
    Let $d,n\geq2$, $\mq\geq-1$, $k\in\NN_0$, $0\leq s<r< k+\frac{1}{2}$, and let $\{\theta_j^*\}_{j=1}^n\subset\SS^d$ have antipodal separation $\underline h>0$. There exists a constant $c_*>0$ such that whenever $2^\mq\underline h\leq c_*$, every $f_n\in L_n^k$ satisfies
    \begin{equation}\label{eqn_inver_highfre}
        \|\pc_\mq f_n\|_{\mathcal{H}^r(\SS^d)}\lesssim 
        \underline{h}^{-(r-s)}\|\pc_\mq f_n\|_{\mathcal{H}^s(\SS^d)}.
    \end{equation}
    Moreover, let $\beta\geq0$ and $\mathfrak L_\beta$ satisfy \eqref{eqn:def_operator}--\eqref{eqn:elliptic_symbol}. For $0\leq s<r< k+\frac{1}{2}-\beta$, and any $v_n\in V_{n,\beta}=\mathfrak{L}_\beta L_n^k(\Theta_n)$,
\begin{equation}\label{eqn:filtered_residual_bernstein}
\|\mathcal P_{\mathfrak q}^c v_n\|_{\mathcal H^{r}(\SS^d)}\lesssim\underline h^{-(r-s)}\|\mathcal P_{\mathfrak q}^c v_n\|_{\mathcal H^{s}(\SS^d)}.
\end{equation}
\end{theorem}

\begin{equation}
\|\mathcal P_{\mathfrak q}^c v_n\|_{\mathcal H^{r}(\SS^d)}\lesssim\underline h^{-(r-s)}\|\mathcal P_{\mathfrak q}^c v_n\|_{\mathcal H^{s}(\SS^d)}.
\end{equation}

\begin{theorem}\label{thm:leastsquare}
Let $d,n\geq2$, $k\in\NN_0$, $\beta\geq0$, $k> \beta+\frac{d-1}{2}$, $r\leq\frac{d+2k+1}{2}-\beta$, $0\leq s<\min\{r,k+\frac{1}{2}-\beta\}$, $\Theta_n=\{\theta_{j,n}^*\}_{j=1}^n\subset\SS^d$, and let $\mathfrak L_\beta$ satisfy \eqref{eqn:def_operator}--\eqref{eqn:elliptic_symbol}. Then there exists a constant $C_2$ such that, whenever $m\geq C_2\underline h^{-d}$, $X=\{\eta_i^*\}_{i=1}^m$ is quasi-uniform, and $f$ satisfies \eqref{eqn:f_evenodd_relu}, the solution $u_{n,m}$ of \eqref{eqn_def_collocation} satisfies
\begin{equation}\label{eqn_main_ls_determ}
\|u-u_{n,m}\|_{\mathcal H^{s+\beta}(\SS^d)}\eqsim\|f-\mathfrak L_\beta u_{n,m}\|_{\mathcal H^s(\SS^d)}\lesssim\underline h^{-s}h^r
\begin{cases}
\|f\|_{\mathcal W^{r,p}(\SS^d)},&\frac dp<r\leq \frac{d}{2},\\
\|f\|_{\mathcal H^r(\SS^d)},&r>\frac{d}{2}.
\end{cases}
\end{equation}
The implied constant is independent of $n,m$, $\Theta_n$, and $X$. If $\{\Theta_n\}_{n\geq1}$ is antipodally quasi-uniform and $m=\lceil C_2\underline h^{-d}\rceil$, then
\begin{equation}
\|u-u_{n,m}\|_{\mathcal H^{s+\beta}(\SS^d)}\eqsim\|f-\mathfrak L_\beta u_{n,m}\|_{\mathcal H^s(\SS^d)}\lesssim n^{-\frac{r-s}{d}}
\begin{cases}
\|f\|_{\mathcal W^{r,p}(\SS^d)},&\frac dp<r\leq \frac{d}{2},\\
\|f\|_{\mathcal H^r(\SS^d)},&r>\frac{d}{2}.
\end{cases}
\end{equation}
\end{theorem}

\begin{lemma}[Approximation by polynomials]\label{lem_sharp_appr_poly}
    Let $r>0$, $0\leq s\leq r$, $1\leq p\leq\infty$. Then for $\mq\in\NN_0$,
    \begin{equation}\label{eqn_sharp_appr_poly}
        \|u-\mP_\mq u\|_{\mathcal W^{s,p}(\SS^d)}\lesssim2^{-\mq (r-s)}\|u\|_{\mathcal{W}^{r,p}(\SS^d)}.
    \end{equation}
\end{lemma}

\begin{lemma}[$\mathcal L^p$ discrete norm estimate]\label{lem_discrete_by_cont}
Let $2\leq p\leq\infty$ and let $X=\{\eta_i^*\}_{i=1}^m\subset\SS^d$ be quasi-uniform with fill distance $h_X$. If
\begin{equation*}
    \frac dp<\tau<\frac dp+1,
\end{equation*}
then every $v\in\mathcal W^{\tau,p}(\SS^d)$ satisfies
\begin{equation}\label{eqn:norming}
    \|v\|_{\ell^2(X)}^2\lesssim\|v\|_{\mathcal L^p(\SS^d)}^2+h_X^{2\tau}\|v\|_{\mathcal W^{\tau,p}(\SS^d)}^2.
\end{equation}
The implied constant depends only on $d$, $p$, and $\tau$.
\end{lemma}

\begin{lemma}[Marcinkiewicz–Zygmund inequality for polynomials]\label{lem_MZ_L2}
Let $X=\{\eta_i^*\}_{i=1}^m$ be quasi-uniform and fill distance $h_X$. For a sufficiently small constant $c_2>0$ and every polynomial $v$ with $\deg(v)\leq c_2h_X^{-1}$,
    \begin{equation}
        \|v\|_{\ell^2(X)}\eqsim\|v\|_{\mathcal{L}^2(\SS^d)}.
    \end{equation}
\end{lemma}

\begin{theorem}\label{thm:main_sphere}
Let $d,n\in\NN$, $k\in\NN_0$, $r\leq\frac{d+2k+1}{2}$, and $\{\theta_j^*\}_{j=1}^n\subset\SS^d$. For any $u\in\mathcal H^r(\SS^d)$ satisfying the parity in \eqref{eqn:f_evenodd_relu}, there exists $u_n\in L_n^k$ such that
\begin{equation}
\|u-u_n\|_{\mathcal H^s(\SS^d)}\lesssim h^{r-s}\|u\|_{\mathcal H^r(\SS^d)},
\end{equation}
holds for $0\leq s<\min\{k+\frac12,r\}$ and
\begin{equation}\label{eqn:trunc_equal}
\widehat u_n(\nu,\ell)=\widehat u(\nu,\ell),\qquad \nu\leq C_1 h^{-1},
\end{equation}
where $C_1$ is a constant independent of $n$ and $\{\theta_j^*\}_{j=1}^n$.
\end{theorem}

\begin{proof}
Let $u=\mathfrak L_\beta^{-1}f$, and let $u_n$ be the function in Theorem~\ref{thm:main_sphere} obtained by applying that theorem to $u$ with smoothness parameters $r+\beta$ and $s+\beta$. Since $\mathfrak L_\beta$ is an elliptic spectral multiplier of order $\beta$,
\begin{equation}\label{eqn_net_appr_u}
    \|\mathfrak L_\beta u_n-f\|_{\mathcal H^s(\SS^d)}
    \eqsim\|u_n-u\|_{\mathcal H^{s+\beta}(\SS^d)}
    \lesssim h^{r-s}\|u\|_{\mathcal H^{r+\beta}(\SS^d)}
    \eqsim h^{r-s}\|f\|_{\mathcal H^r(\SS^d)}.
\end{equation}
Here $\mathfrak L_\beta^{-1}$ preserves spherical harmonic degrees and hence preserves the parity condition required in Theorem~\ref{thm:main_sphere}. Let
\begin{equation*}
    g=\mathfrak L_\beta(u_{n,m}-u_n)\in V_{n,\beta}.
\end{equation*}
We first estimate $\|g\|_{\mathcal L^2(\SS^d)}$. Choose
$$
\frac d2<\tau<\min\big\{k+\frac12-\beta,\frac{d}{2}+1\big\},
$$
and set $\mathfrak p=\lfloor\log_2(c_2m^{1/d})\rfloor-1$, where $c_2$ is the constant in Lemma~\ref{lem_MZ_L2}. Then
\begin{equation}\label{eqn_samp_step1}
    \begin{split}
        \|g\|_{\mathcal L^2(\SS^d)}
        \leq&\|\mathcal P_{\mathfrak p}g\|_{\mathcal L^2(\SS^d)}
        +\|g-\mathcal P_{\mathfrak p}g\|_{\mathcal L^2(\SS^d)}\\
        \lesssim&\|\mathcal P_{\mathfrak p}g\|_{\ell^2(X)}
        +m^{-\frac{\tau}{d}}\|g\|_{\mathcal H^\tau(\SS^d)},
    \end{split}
\end{equation}
where the last inequality follows by applying Lemma~\ref{lem_MZ_L2} to $\mathcal P_{\mathfrak p}g$ and Lemma~\ref{lem_sharp_appr_poly} to $g$.

Applying Lemma~\ref{lem_discrete_by_cont} to $\pc_{\mathfrak p}g$ gives
\begin{equation}\label{eqn_est_Pc_ell2}
    \begin{split}
        \|\pc_{\mathfrak p}g\|_{\ell^2(X)}^2
        \lesssim&\|\pc_{\mathfrak p}g\|_{\mathcal L^2(\SS^d)}^2
        +m^{-\frac{2\tau}{d}}\|\pc_{\mathfrak p}g\|_{\mathcal H^\tau(\SS^d)}^2\\
        \lesssim&\left(2^{-2\mathfrak p\tau}+m^{-\frac{2\tau}{d}}\right)
        \|g\|_{\mathcal H^\tau(\SS^d)}^2
        \lesssim m^{-\frac{2\tau}{d}}\|g\|_{\mathcal H^\tau(\SS^d)}^2,
    \end{split}
\end{equation}
where the second inequality follows from Lemma~\ref{lem_sharp_appr_poly} and the boundedness of $\pc_{\mathfrak p}$ on $\mathcal H^\tau(\SS^d)$. Combining \eqref{eqn_samp_step1} and \eqref{eqn_est_Pc_ell2}, we obtain
\begin{equation}\label{eqn:L2_by_discrete}
    \begin{split}
        \|g\|_{\mathcal L^2(\SS^d)}
        \lesssim&\left(\|\pc_{\mathfrak p}g\|_{\ell^2(X)}+\|g\|_{\ell^2(X)}\right)
        +m^{-\frac{\tau}{d}}\|g\|_{\mathcal H^\tau(\SS^d)}\\
        \lesssim&\|g\|_{\ell^2(X)}+m^{-\frac{\tau}{d}}\|g\|_{\mathcal H^\tau(\SS^d)}.
    \end{split}
\end{equation}

\medskip
Since $g\in V_{n,\beta}$, the residual Bernstein estimate \eqref{eqn:residual_bernstein} from Theorem~\ref{thm:inverse} gives
\begin{equation}\label{eqn:residual_bernstein_used}
    \|g\|_{\mathcal H^\tau(\SS^d)}
    \lesssim\underline h^{-\tau}\|g\|_{\mathcal L^2(\SS^d)},
    \qquad
    \|g\|_{\mathcal H^s(\SS^d)}
    \lesssim\underline h^{-s}\|g\|_{\mathcal L^2(\SS^d)}.
\end{equation}
For $s=0$, the second inequality is understood trivially. It follows from \eqref{eqn:L2_by_discrete} that
\begin{equation*}
    \|g\|_{\mathcal L^2(\SS^d)}
    \lesssim\|g\|_{\ell^2(X)}
    +(m\underline h^d)^{-\frac{\tau}{d}}\|g\|_{\mathcal L^2(\SS^d)}.
\end{equation*}
Hence, there exists a constant $C_2>0$ such that, whenever $m\geq C_2\underline h^{-d}$, the last term can be absorbed into the left-hand side. Therefore,
\begin{equation}\label{eqn:residual_stability}
    \|g\|_{\mathcal H^s(\SS^d)}
    \lesssim\underline h^{-s}\|g\|_{\mathcal L^2(\SS^d)}
    \lesssim\underline h^{-s}\|g\|_{\ell^2(X)}.
\end{equation}
Since $u_{n,m}$ is the $\ell^2$-minimizer of $f-\mathfrak L_\beta v_n$ over $L_n^k$, we obtain
\begin{equation}\label{eqn:error_decom2}
    \begin{split}
        \|g\|_{\mathcal H^s(\SS^d)}
        \lesssim&\underline h^{-s}\|\mathfrak L_\beta(u_{n,m}-u_n)\|_{\ell^2(X)}\\
        \lesssim&\underline h^{-s}\|\mathfrak L_\beta u_{n,m}-f\|_{\ell^2(X)}
        +\underline h^{-s}\|f-\mathfrak L_\beta u_n\|_{\ell^2(X)}\\
        \leq&2\underline h^{-s}\|f-\mathfrak L_\beta u_n\|_{\ell^2(X)}.
    \end{split}
\end{equation}

\medskip
We now estimate $\|f-\mathfrak L_\beta u_n\|_{\ell^2(X)}$. Let $C_1$ be the constant in Theorem \ref{thm:main_sphere} and $c_*$ be the constant in Theorem \ref{thm_inverse_hig_freq}. Choose a fixed integer $q_0\geq1$ sufficiently large and set
\begin{equation*}
    \mq=\lfloor\log_2C_1 h^{-1}\rfloor-q_0,
\end{equation*}
so that
\begin{equation*}
    2^{\mq+1}\leq C_1 h^{-1},\qquad 2^\mq\underline h\leq c_*,\qquad 2^{-\mq}\eqsim h.
\end{equation*}
The finitely many coarse cases for which $\mq<0$ can be included by enlarging the implied constant; hence, in the remainder of the proof, we assume $\mq\geq0$. By Theorem~\ref{thm:main_sphere}, $\mathcal P_\mq u=\mathcal P_\mq u_n$. Since $\mathfrak L_\beta$ commutes with $\mathcal P_\mq$, it follows that
\begin{equation*}
    \mathcal P_\mq f=\mathcal P_\mq\mathfrak L_\beta u_n.
\end{equation*}
Consequently,
\begin{equation}\label{eqn:error_decom3}
    \begin{split}
        \|f-\mathfrak L_\beta u_n\|_{\ell^2(X)}
        \leq&\|f-\mathcal P_\mq f\|_{\ell^2(X)}
        +\|\mathcal P_\mq f-\mathfrak L_\beta u_n\|_{\ell^2(X)}\\
        =&\|\pc_\mq f\|_{\ell^2(X)}
        +\|\pc_\mq\mathfrak L_\beta u_n\|_{\ell^2(X)}.
    \end{split}
\end{equation}
The two terms are estimated separately. For $\|\pc_\mq f\|_{\ell^2(X)}$:
\begin{itemize}
    \item If $r\leq d/2$, let $2<p<\infty$ satisfy $r>d/p$ and choose
\begin{equation*}
    \frac{d}{p}<\tau<\min\left\{r,\frac{d}{p}+1\right\}.
\end{equation*}
Applying Lemma~\ref{lem_discrete_by_cont} and Lemma~\ref{lem_sharp_appr_poly} with $p$ yields
\begin{equation*}
    \begin{split}
        \|\pc_\mq f\|_{\ell^2(X)}
        &\lesssim\|\pc_\mq f\|_{\mathcal L^p(\SS^d)}+h_X^\tau\|\pc_\mq f\|_{\mathcal W^{\tau,p}(\SS^d)}\\
        &\lesssim\left(2^{-\mq r}+h_X^\tau2^{-\mq(r-\tau)}\right)\|f\|_{\mathcal W^{r,p}(\SS^d)}
        \lesssim h^r\|f\|_{\mathcal W^{r,p}(\SS^d)}.
    \end{split}
\end{equation*}
    \item If $r>d/2$, choose
    \begin{equation*}
        \frac d2<\tau_1<\min\left\{r,\frac d2+1\right\}.
    \end{equation*}
    Applying Lemmas~\ref{lem_discrete_by_cont} and \ref{lem_sharp_appr_poly}, and using $m^{-1/d}\lesssim\underline h\lesssim h$, we obtain
    \begin{equation*}
        \begin{split}
            \|\pc_\mq f\|_{\ell^2(X)}
            \lesssim&\|\pc_\mq f\|_{\mathcal L^2(\SS^d)}
            +m^{-\frac{\tau_1}{d}}\|\pc_\mq f\|_{\mathcal H^{\tau_1}(\SS^d)}\\
            \lesssim&\left(h^r+m^{-\frac{\tau_1}{d}}h^{r-\tau_1}\right)
            \|f\|_{\mathcal H^r(\SS^d)}
            \lesssim h^r\|f\|_{\mathcal H^r(\SS^d)}.
        \end{split}
    \end{equation*}
\end{itemize}
Therefore,
\begin{equation}\label{eqn_poincare1}
    \|\pc_\mq f\|_{\ell^2(X)}
    \lesssim h^r
    \left\{\begin{array}{ll}
        \|f\|_{\mathcal W^{r,p}(\SS^d)},&\frac{d}{p}<r\leq\frac d2,\\[0.3em]
        \|f\|_{\mathcal H^r(\SS^d)},&r>\frac d2.
    \end{array}\right.
\end{equation}

\medskip
For $\|\pc_\mq\mathfrak L_\beta u_n\|_{\ell^2(X)}$, we apply Lemma~\ref{lem_discrete_by_cont} and the filtered residual Bernstein inequality \eqref{eqn:filtered_residual_bernstein}:
\begin{equation}\label{eqn_poincare2}
    \begin{split}
        \|\pc_\mq\mathfrak L_\beta u_n\|_{\ell^2(X)}
        \lesssim&\|\pc_\mq\mathfrak L_\beta u_n\|_{\mathcal L^2(\SS^d)}
        +m^{-\frac{\tau}{d}}\|\pc_\mq\mathfrak L_\beta u_n\|_{\mathcal H^\tau(\SS^d)}\\
        \lesssim&\|\pc_\mq\mathfrak L_\beta u_n\|_{\mathcal L^2(\SS^d)}
        +m^{-\frac{\tau}{d}}\underline h^{-\tau}
        \|\pc_\mq\mathfrak L_\beta u_n\|_{\mathcal L^2(\SS^d)}\\
        \lesssim&\|\pc_\mq\mathfrak L_\beta u_n\|_{\mathcal L^2(\SS^d)}=\|\pc_\mq f+(\mathfrak L_\beta u_n-f)\|_{\mathcal L^2(\SS^d)}
        \lesssim h^r\|f\|_{\mathcal H^r(\SS^d)}.
    \end{split}
\end{equation}
Here the second inequality follows from \eqref{eqn:filtered_residual_bernstein}, the third one uses $m\geq C_2\underline h^{-d}$, and the equality in the last line follows from
\begin{equation*}
    \mathcal P_\mq(f-\mathfrak L_\beta u_n)=0.
\end{equation*}
The last inequality follows from Lemma~\ref{lem_sharp_appr_poly} and \eqref{eqn_net_appr_u} with $s=0$. In the case $r\leq d/2$, we also use
\begin{equation*}
    \|f\|_{\mathcal H^r(\SS^d)}\lesssim\|f\|_{\mathcal W^{r,p}(\SS^d)},\quad p\geq2.
\end{equation*}

\medskip
Combining \eqref{eqn:error_decom2}--\eqref{eqn_poincare2}, we obtain
\begin{equation}
    \begin{split}
        \|g\|_{\mathcal H^s(\SS^d)}
        \lesssim&\underline h^{-s}\left(\|\pc_\mq\mathfrak L_\beta u_n\|_{\ell^2(X)}
        +\|\pc_\mq f\|_{\ell^2(X)}\right)\\
        \lesssim&\underline h^{-s}h^r
        \left\{\begin{array}{ll}
            \|f\|_{\mathcal W^{r,p}(\SS^d)},&\frac{d}{2}<r\leq\frac d2,\\[0.3em]
            \|f\|_{\mathcal H^r(\SS^d)},&r>\frac d2.
        \end{array}\right.
    \end{split}
\end{equation}

Together with \eqref{eqn_net_appr_u}, this gives
\begin{equation}
    \begin{split}
        \|\mathfrak L_\beta u_{n,m}-f\|_{\mathcal H^s(\SS^d)}
        \leq&\|g\|_{\mathcal H^s(\SS^d)}
        +\|f-\mathfrak L_\beta u_n\|_{\mathcal H^s(\SS^d)}\\
        \lesssim&\underline h^{-s}h^r
        \left\{\begin{array}{ll}
            \|f\|_{\mathcal W^{r,p}(\SS^d)},&\frac dp<r\leq\frac d2,\\[0.3em]
            \|f\|_{\mathcal H^r(\SS^d)},&r>\frac d2.
        \end{array}\right.,
    \end{split}
\end{equation}
where we have used $\underline h\lesssim h$ to absorb the term $h^{r-s}$. This proves Theorem~\ref{thm:leastsquare}.
\end{proof}

\end{document}
